\section{Conclusion}
\label{sec:conclusion}

A model asked about exercise by a user who mentions a disability answers a different
question, one about risk and supervision. Removing that behaviour at inference time is a
reasonable goal, and the field measures progress toward it with a metric that counts the
clinical vocabulary the model stops producing.

That metric cannot tell the difference between a model that stopped medicalizing and a
model that stopped writing. We found three methods sitting at settings where it reports
success and the responses are corrupted, looped, or truncated, and the failure modes
differ enough that no single filter on emptiness, length or uniqueness catches them all.
The method with the lowest score in our comparison is the one that damages the output
most.

Measuring quality alongside the target metric changes the ranking and leaves a usable
picture: a narrow band of configurations that reduce medicalization while leaving the
text intact, occupied almost entirely by multi-rank projection at the operating point its
rank prescribes. That rank matters is a separate and simpler claim than the one we set
out to make, and it survives the control experiments that the stronger claim does not.

The practical recommendation is small and cheap. Report a quality measure beside the
target metric, since a vocabulary-absence score is satisfied by absence of any kind. And
report the number of independent base prompts alongside the number of contrastive pairs,
since the spectral statistics read off those pairs move further under a change of prompt
diversity than under a change of bias.
